% --- Archon physics formalization source begin ---
% archon:physics
% archon:covers PhyXMiniProblems/problem_phyx_mini_0706.lean
% archon:source-report reports/phyx_mini/problem_phyx_mini_0706.source.json

\chapter{Physics problem phyx\_mini\_0706}
\label{ch:PhyXMiniProblems_problem_phyx_mini_0706}

\paragraph{Problem source.}
\#\# Physical scenario

A \textbackslash{}(2.0\textbackslash{},\textbackslash{}mathrm\{kg\}\textbackslash{}) bucket is attached to a massless string that is wrapped around a \textbackslash{}(1.0\textbackslash{},\textbackslash{}mathrm\{kg\}\textbackslash{}), \textbackslash{}(4.0\textbackslash{},\textbackslash{}mathrm\{cm\}\textbackslash{})-diameter cylinder, as shown in figure. The cylinder rotates on an axle through the center. The bucket is released from rest \textbackslash{}(1.0\textbackslash{},\textbackslash{}mathrm\{m\}\textbackslash{}) above the floor.

\#\# Question

How long does it take to reach the floor?

\#\# Answer choices

- A: A: 0.55s

- B: B: 0.45s

- C: C: 0.50s

- D: D: 0.40s

\#\# Figure caption (auxiliary; use the image as primary evidence)

The image depicts a physics problem involving a pulley system. Here's a detailed description:

1. **Objects and Setup:**
   - A cylindrical pulley with radius \textbackslash{}( R \textbackslash{}) and mass \textbackslash{}( M \textbackslash{}) is shown at the top, mounted on an axle.
   - A bucket is attached to the pulley via a rope. The bucket is positioned vertically below the pulley.

2. **Pulley:**
   - Marked with a radius \textbackslash{}( R = 2.0 \textbackslash{}, \textbackslash{}text\{cm\} \textbackslash{}) and mass \textbackslash{}( M = 1.0 \textbackslash{}, \textbackslash{}text\{kg\} \textbackslash{}).
   - An arrow indicates rotational motion of the pulley, marked as \textbackslash{}( \textbackslash{}alpha \textbackslash{}).

3. **Bucket:**
   - A bucket with mass \textbackslash{}( m = 2.0 \textbackslash{}, \textbackslash{}text\{kg\} \textbackslash{}) is shown.
   - It’s initially positioned at \textbackslash{}( y\_0 = 1.0 \textbackslash{}, \textbackslash{}text\{m\} \textbackslash{}) above the ground.
   - The initial velocity of the bucket is \textbackslash{}( v\_0 = 0 \textbackslash{}, \textbackslash{}text\{m/s\} \textbackslash{}).
   - The bucket is accelerating downward, denoted by an arrow labeled \textbackslash{}( a \textbackslash{}).

4. **Vertical Positions:**
   - The setup shows two vertical positions marked on the side:
     - Initial position \textbackslash{}( y\_0 = 1.0 \textbackslash{}, \textbackslash{}text\{m\} \textbackslash{}).
     - Final position \textbackslash{}( y\_1 = 0 \textbackslash{}, \textbackslash{}text\{m\} \textbackslash{}) (ground level).

5. **Coordinate System:**
   - A vertical axis labeled \textbackslash{}( y \textbackslash{}) is present, indicating the direction and positions.

This setup likely represents a classic mechanics problem investigating rotational motion and linear acceleration.

\#\# Dataset metadata

Rotational Motion; Physical Model Grounding Reasoning, Multi-Formula Reasoning

\paragraph{Recorded answer/context.}
C: C: 0.50s

\paragraph{Figure/image path.}
/root/proposal\_for\_physic/science-mango/phyx\_mini\_run/phyx\_data/test\_image/706.png

\paragraph{Formalization target.}
create a compiling Lean file with sorry bodies at `PhyXMiniProblems/problem\_phyx\_mini\_0706.lean`. The Lean declarations must preserve the physical quantities, dimensions or dimensional roles, figure labels, governing-law hypotheses, and final relation expressed by this problem.
Use Mathlib/Physlib names found through LeanExplore where available. If a physics API is missing, introduce faithful local abstractions rather than scalar placeholder aliases.

\begin{theorem}[Physics formalization target]
\label{thm:physics:phyx_mini_0706:target}
\lean{PhyXMiniProblems.ProblemPhyXMini0706.problem_phyx_mini_0706}
\uses{def:physics:phyx-mini-0706:phyxminiproblems-problemphyxmini0706-massinkilograms, def:physics:phyx-mini-0706:phyxminiproblems-problemphyxmini0706-signedlengthinmeters, def:physics:phyx-mini-0706:phyxminiproblems-problemphyxmini0706-timeinseconds, def:physics:phyx-mini-0706:phyxminiproblems-problemphyxmini0706-accelerationinmeterspersecondsquared, def:physics:phyx-mini-0706:phyxminiproblems-problemphyxmini0706-bucketcylindersetup, def:physics:phyx-mini-0706:phyxminiproblems-problemphyxmini0706-matchesproblemandfigurereadouts, def:physics:phyx-mini-0706:phyxminiproblems-problemphyxmini0706-usesstandardearthgravity, def:physics:phyx-mini-0706:phyxminiproblems-problemphyxmini0706-hasphysicalparameters, def:physics:phyx-mini-0706:phyxminiproblems-problemphyxmini0706-satisfiescylindergeometry, def:physics:phyx-mini-0706:phyxminiproblems-problemphyxmini0706-satisfiesuniformsolidcylinderinertialaw, def:physics:phyx-mini-0706:phyxminiproblems-problemphyxmini0706-satisfiesbucketcylinderdynamics, def:physics:phyx-mini-0706:phyxminiproblems-problemphyxmini0706-satisfiesconstantaccelerationkinematics, def:physics:phyx-mini-0706:phyxminiproblems-problemphyxmini0706-isuniqueclosestdisplayedtime}
The assigned autoformalize agent should translate this physics problem into Lean declarations in the covered file, with theorem and lemma proof bodies written as `by sorry`.
\end{theorem}
\begin{proof}
This is an autoformalization task, not a proof task. Produce statements that can later be proved by the physics prover without weakening the physical model.
\end{proof}

% --- Archon declaration topology begin ---
\section{Declaration topology}
\begin{definition}[Mass Quantity]
  \label{def:physics:phyx-mini-0706:phyxminiproblems-problemphyxmini0706-massquantity}
  \lean{PhyXMiniProblems.ProblemPhyXMini0706.MassQuantity}
  A nonnegative physical mass
\end{definition}
\begin{proof}
  This is the declared model definition of Mass Quantity.
\end{proof}
\begin{definition}[Length Quantity]
  \label{def:physics:phyx-mini-0706:phyxminiproblems-problemphyxmini0706-lengthquantity}
  \lean{PhyXMiniProblems.ProblemPhyXMini0706.LengthQuantity}
  A nonnegative physical length
\end{definition}
\begin{proof}
  This is the declared model definition of Length Quantity.
\end{proof}
\begin{definition}[Signed Length Quantity]
  \label{def:physics:phyx-mini-0706:phyxminiproblems-problemphyxmini0706-signedlengthquantity}
  \lean{PhyXMiniProblems.ProblemPhyXMini0706.SignedLengthQuantity}
  A signed vertical coordinate
\end{definition}
\begin{proof}
  This is the declared model definition of Signed Length Quantity.
\end{proof}
\begin{definition}[Time Quantity]
  \label{def:physics:phyx-mini-0706:phyxminiproblems-problemphyxmini0706-timequantity}
  \lean{PhyXMiniProblems.ProblemPhyXMini0706.TimeQuantity}
  A nonnegative physical duration
\end{definition}
\begin{proof}
  This is the declared model definition of Time Quantity.
\end{proof}
\begin{definition}[Signed Speed Quantity]
  \label{def:physics:phyx-mini-0706:phyxminiproblems-problemphyxmini0706-signedspeedquantity}
  \lean{PhyXMiniProblems.ProblemPhyXMini0706.SignedSpeedQuantity}
  A signed linear velocity component
\end{definition}
\begin{proof}
  This is the declared model definition of Signed Speed Quantity.
\end{proof}
\begin{definition}[Acceleration Quantity]
  \label{def:physics:phyx-mini-0706:phyxminiproblems-problemphyxmini0706-accelerationquantity}
  \lean{PhyXMiniProblems.ProblemPhyXMini0706.AccelerationQuantity}
  A nonnegative linear-acceleration magnitude
\end{definition}
\begin{proof}
  This is the declared model definition of Acceleration Quantity.
\end{proof}
\begin{definition}[Force Quantity]
  \label{def:physics:phyx-mini-0706:phyxminiproblems-problemphyxmini0706-forcequantity}
  \lean{PhyXMiniProblems.ProblemPhyXMini0706.ForceQuantity}
  A nonnegative force magnitude
\end{definition}
\begin{proof}
  This is the declared model definition of Force Quantity.
\end{proof}
\begin{definition}[Moment Of Inertia Quantity]
  \label{def:physics:phyx-mini-0706:phyxminiproblems-problemphyxmini0706-momentofinertiaquantity}
  \lean{PhyXMiniProblems.ProblemPhyXMini0706.MomentOfInertiaQuantity}
  Moment of inertia has dimension mass times length squared
\end{definition}
\begin{proof}
  This is the declared model definition of Moment Of Inertia Quantity.
\end{proof}
\begin{definition}[Angular Acceleration Quantity]
  \label{def:physics:phyx-mini-0706:phyxminiproblems-problemphyxmini0706-angularaccelerationquantity}
  \lean{PhyXMiniProblems.ProblemPhyXMini0706.AngularAccelerationQuantity}
  Angular acceleration has inverse-time-squared dimension because radians are dimensionless
\end{definition}
\begin{proof}
  This is the declared model definition of Angular Acceleration Quantity.
\end{proof}
\begin{definition}[mass Readout]
  \label{def:physics:phyx-mini-0706:phyxminiproblems-problemphyxmini0706-massreadout}
  \lean{PhyXMiniProblems.ProblemPhyXMini0706.massReadout}
  \uses{def:physics:phyx-mini-0706:phyxminiproblems-problemphyxmini0706-massquantity}
  Read a physical mass in a selected mass unit
\end{definition}
\begin{proof}
  This is the declared model definition of mass Readout.
\end{proof}
\begin{definition}[length Readout]
  \label{def:physics:phyx-mini-0706:phyxminiproblems-problemphyxmini0706-lengthreadout}
  \lean{PhyXMiniProblems.ProblemPhyXMini0706.lengthReadout}
  \uses{def:physics:phyx-mini-0706:phyxminiproblems-problemphyxmini0706-lengthquantity}
  Read a nonnegative length in a selected length unit
\end{definition}
\begin{proof}
  This is the declared model definition of length Readout.
\end{proof}
\begin{definition}[signed Length Readout]
  \label{def:physics:phyx-mini-0706:phyxminiproblems-problemphyxmini0706-signedlengthreadout}
  \lean{PhyXMiniProblems.ProblemPhyXMini0706.signedLengthReadout}
  \uses{def:physics:phyx-mini-0706:phyxminiproblems-problemphyxmini0706-signedlengthquantity}
  Read a signed coordinate in a selected length unit
\end{definition}
\begin{proof}
  This is the declared model definition of signed Length Readout.
\end{proof}
\begin{definition}[time Readout]
  \label{def:physics:phyx-mini-0706:phyxminiproblems-problemphyxmini0706-timereadout}
  \lean{PhyXMiniProblems.ProblemPhyXMini0706.timeReadout}
  \uses{def:physics:phyx-mini-0706:phyxminiproblems-problemphyxmini0706-timequantity}
  Read a duration in a selected time unit
\end{definition}
\begin{proof}
  This is the declared model definition of time Readout.
\end{proof}
\begin{definition}[signed Speed Readout]
  \label{def:physics:phyx-mini-0706:phyxminiproblems-problemphyxmini0706-signedspeedreadout}
  \lean{PhyXMiniProblems.ProblemPhyXMini0706.signedSpeedReadout}
  \uses{def:physics:phyx-mini-0706:phyxminiproblems-problemphyxmini0706-signedspeedquantity}
  Read a signed speed in coherent selected length and time units
\end{definition}
\begin{proof}
  This is the declared model definition of signed Speed Readout.
\end{proof}
\begin{definition}[acceleration Readout]
  \label{def:physics:phyx-mini-0706:phyxminiproblems-problemphyxmini0706-accelerationreadout}
  \lean{PhyXMiniProblems.ProblemPhyXMini0706.accelerationReadout}
  \uses{def:physics:phyx-mini-0706:phyxminiproblems-problemphyxmini0706-accelerationquantity}
  Read an acceleration magnitude in coherent selected units
\end{definition}
\begin{proof}
  This is the declared model definition of acceleration Readout.
\end{proof}
\begin{definition}[force Readout]
  \label{def:physics:phyx-mini-0706:phyxminiproblems-problemphyxmini0706-forcereadout}
  \lean{PhyXMiniProblems.ProblemPhyXMini0706.forceReadout}
  \uses{def:physics:phyx-mini-0706:phyxminiproblems-problemphyxmini0706-forcequantity}
  Read a force magnitude in coherent selected mechanical units
\end{definition}
\begin{proof}
  This is the declared model definition of force Readout.
\end{proof}
\begin{definition}[moment Of Inertia Readout]
  \label{def:physics:phyx-mini-0706:phyxminiproblems-problemphyxmini0706-momentofinertiareadout}
  \lean{PhyXMiniProblems.ProblemPhyXMini0706.momentOfInertiaReadout}
  \uses{def:physics:phyx-mini-0706:phyxminiproblems-problemphyxmini0706-momentofinertiaquantity}
  Read a moment of inertia in coherent selected mass and length units
\end{definition}
\begin{proof}
  This is the declared model definition of moment Of Inertia Readout.
\end{proof}
\begin{definition}[angular Acceleration Readout]
  \label{def:physics:phyx-mini-0706:phyxminiproblems-problemphyxmini0706-angularaccelerationreadout}
  \lean{PhyXMiniProblems.ProblemPhyXMini0706.angularAccelerationReadout}
  \uses{def:physics:phyx-mini-0706:phyxminiproblems-problemphyxmini0706-angularaccelerationquantity}
  Read angular acceleration in inverse square units of the selected time unit
\end{definition}
\begin{proof}
  This is the declared model definition of angular Acceleration Readout.
\end{proof}
\begin{definition}[mass In Kilograms]
  \label{def:physics:phyx-mini-0706:phyxminiproblems-problemphyxmini0706-massinkilograms}
  \lean{PhyXMiniProblems.ProblemPhyXMini0706.massInKilograms}
  \uses{def:physics:phyx-mini-0706:phyxminiproblems-problemphyxmini0706-massquantity, def:physics:phyx-mini-0706:phyxminiproblems-problemphyxmini0706-massreadout}
  SI kilogram readout of a mass
\end{definition}
\begin{proof}
  This is the declared model definition of mass In Kilograms.
\end{proof}
\begin{definition}[length In Meters]
  \label{def:physics:phyx-mini-0706:phyxminiproblems-problemphyxmini0706-lengthinmeters}
  \lean{PhyXMiniProblems.ProblemPhyXMini0706.lengthInMeters}
  \uses{def:physics:phyx-mini-0706:phyxminiproblems-problemphyxmini0706-lengthquantity, def:physics:phyx-mini-0706:phyxminiproblems-problemphyxmini0706-lengthreadout}
  SI meter readout of a nonnegative length
\end{definition}
\begin{proof}
  This is the declared model definition of length In Meters.
\end{proof}
\begin{definition}[signed Length In Meters]
  \label{def:physics:phyx-mini-0706:phyxminiproblems-problemphyxmini0706-signedlengthinmeters}
  \lean{PhyXMiniProblems.ProblemPhyXMini0706.signedLengthInMeters}
  \uses{def:physics:phyx-mini-0706:phyxminiproblems-problemphyxmini0706-signedlengthquantity, def:physics:phyx-mini-0706:phyxminiproblems-problemphyxmini0706-signedlengthreadout}
  SI meter readout of a signed vertical coordinate
\end{definition}
\begin{proof}
  This is the declared model definition of signed Length In Meters.
\end{proof}
\begin{definition}[time In Seconds]
  \label{def:physics:phyx-mini-0706:phyxminiproblems-problemphyxmini0706-timeinseconds}
  \lean{PhyXMiniProblems.ProblemPhyXMini0706.timeInSeconds}
  \uses{def:physics:phyx-mini-0706:phyxminiproblems-problemphyxmini0706-timequantity, def:physics:phyx-mini-0706:phyxminiproblems-problemphyxmini0706-timereadout}
  SI second readout of a duration
\end{definition}
\begin{proof}
  This is the declared model definition of time In Seconds.
\end{proof}
\begin{definition}[signed Speed In Meters Per Second]
  \label{def:physics:phyx-mini-0706:phyxminiproblems-problemphyxmini0706-signedspeedinmeterspersecond}
  \lean{PhyXMiniProblems.ProblemPhyXMini0706.signedSpeedInMetersPerSecond}
  \uses{def:physics:phyx-mini-0706:phyxminiproblems-problemphyxmini0706-signedspeedquantity, def:physics:phyx-mini-0706:phyxminiproblems-problemphyxmini0706-signedspeedreadout}
  SI meter-per-second readout of a signed velocity
\end{definition}
\begin{proof}
  This is the declared model definition of signed Speed In Meters Per Second.
\end{proof}
\begin{definition}[acceleration In Meters Per Second Squared]
  \label{def:physics:phyx-mini-0706:phyxminiproblems-problemphyxmini0706-accelerationinmeterspersecondsquared}
  \lean{PhyXMiniProblems.ProblemPhyXMini0706.accelerationInMetersPerSecondSquared}
  \uses{def:physics:phyx-mini-0706:phyxminiproblems-problemphyxmini0706-accelerationquantity, def:physics:phyx-mini-0706:phyxminiproblems-problemphyxmini0706-accelerationreadout}
  SI meter-per-second-squared readout of an acceleration magnitude
\end{definition}
\begin{proof}
  This is the declared model definition of acceleration In Meters Per Second Squared.
\end{proof}
\begin{definition}[Cylinder Mass Model]
  \label{def:physics:phyx-mini-0706:phyxminiproblems-problemphyxmini0706-cylindermassmodel}
  \lean{PhyXMiniProblems.ProblemPhyXMini0706.CylinderMassModel}
  How the cylinder's mass is idealized
\end{definition}
\begin{proof}
  This is the declared model definition of Cylinder Mass Model.
\end{proof}
\begin{definition}[String Mass Model]
  \label{def:physics:phyx-mini-0706:phyxminiproblems-problemphyxmini0706-stringmassmodel}
  \lean{PhyXMiniProblems.ProblemPhyXMini0706.StringMassModel}
  The mass model of the connecting string
\end{definition}
\begin{proof}
  This is the declared model definition of String Mass Model.
\end{proof}
\begin{definition}[Axle Location]
  \label{def:physics:phyx-mini-0706:phyxminiproblems-problemphyxmini0706-axlelocation}
  \lean{PhyXMiniProblems.ProblemPhyXMini0706.AxleLocation}
  Where the rotation axle passes through the cylinder
\end{definition}
\begin{proof}
  This is the declared model definition of Axle Location.
\end{proof}
\begin{definition}[Axle Resistance Model]
  \label{def:physics:phyx-mini-0706:phyxminiproblems-problemphyxmini0706-axleresistancemodel}
  \lean{PhyXMiniProblems.ProblemPhyXMini0706.AxleResistanceModel}
  The resistance supplied by the axle
\end{definition}
\begin{proof}
  This is the declared model definition of Axle Resistance Model.
\end{proof}
\begin{definition}[String Cylinder Contact]
  \label{def:physics:phyx-mini-0706:phyxminiproblems-problemphyxmini0706-stringcylindercontact}
  \lean{PhyXMiniProblems.ProblemPhyXMini0706.StringCylinderContact}
  Whether the string slides relative to the cylinder surface
\end{definition}
\begin{proof}
  This is the declared model definition of String Cylinder Contact.
\end{proof}
\begin{definition}[Figure Label]
  \label{def:physics:phyx-mini-0706:phyxminiproblems-problemphyxmini0706-figurelabel}
  \lean{PhyXMiniProblems.ProblemPhyXMini0706.FigureLabel}
  Text or arrow labels visible in the supplied figure
\end{definition}
\begin{proof}
  This is the declared model definition of Figure Label.
\end{proof}
\begin{definition}[Supplied Bucket Cylinder Figure]
  \label{def:physics:phyx-mini-0706:phyxminiproblems-problemphyxmini0706-suppliedbucketcylinderfigure}
  \lean{PhyXMiniProblems.ProblemPhyXMini0706.SuppliedBucketCylinderFigure}
  \uses{def:physics:phyx-mini-0706:phyxminiproblems-problemphyxmini0706-figurelabel}
  Qualitative geometry transcribed from the primary image
\end{definition}
\begin{proof}
  This is the declared model definition of Supplied Bucket Cylinder Figure.
\end{proof}
\begin{definition}[Bucket Cylinder Setup]
  \label{def:physics:phyx-mini-0706:phyxminiproblems-problemphyxmini0706-bucketcylindersetup}
  \lean{PhyXMiniProblems.ProblemPhyXMini0706.BucketCylinderSetup}
  \uses{def:physics:phyx-mini-0706:phyxminiproblems-problemphyxmini0706-massquantity, def:physics:phyx-mini-0706:phyxminiproblems-problemphyxmini0706-lengthquantity, def:physics:phyx-mini-0706:phyxminiproblems-problemphyxmini0706-signedlengthquantity, def:physics:phyx-mini-0706:phyxminiproblems-problemphyxmini0706-timequantity, def:physics:phyx-mini-0706:phyxminiproblems-problemphyxmini0706-signedspeedquantity, def:physics:phyx-mini-0706:phyxminiproblems-problemphyxmini0706-accelerationquantity, def:physics:phyx-mini-0706:phyxminiproblems-problemphyxmini0706-forcequantity, def:physics:phyx-mini-0706:phyxminiproblems-problemphyxmini0706-momentofinertiaquantity, def:physics:phyx-mini-0706:phyxminiproblems-problemphyxmini0706-angularaccelerationquantity, def:physics:phyx-mini-0706:phyxminiproblems-problemphyxmini0706-cylindermassmodel, def:physics:phyx-mini-0706:phyxminiproblems-problemphyxmini0706-stringmassmodel, def:physics:phyx-mini-0706:phyxminiproblems-problemphyxmini0706-axlelocation, def:physics:phyx-mini-0706:phyxminiproblems-problemphyxmini0706-axleresistancemodel, def:physics:phyx-mini-0706:phyxminiproblems-problemphyxmini0706-stringcylindercontact, def:physics:phyx-mini-0706:phyxminiproblems-problemphyxmini0706-suppliedbucketcylinderfigure}
  Independent physical quantities and response fields for the modeled motion. The acceleration, tension, inertia, angular acceleration, and fall time are not defined from the requested answer; the governing-law predicates below constrain them
\end{definition}
\begin{proof}
  This is the declared model definition of Bucket Cylinder Setup.
\end{proof}
\begin{definition}[Matches Problem And Figure Readouts]
  \label{def:physics:phyx-mini-0706:phyxminiproblems-problemphyxmini0706-matchesproblemandfigurereadouts}
  \lean{PhyXMiniProblems.ProblemPhyXMini0706.MatchesProblemAndFigureReadouts}
  \uses{def:physics:phyx-mini-0706:phyxminiproblems-problemphyxmini0706-signedspeedreadout, def:physics:phyx-mini-0706:phyxminiproblems-problemphyxmini0706-massinkilograms, def:physics:phyx-mini-0706:phyxminiproblems-problemphyxmini0706-lengthinmeters, def:physics:phyx-mini-0706:phyxminiproblems-problemphyxmini0706-signedlengthinmeters, def:physics:phyx-mini-0706:phyxminiproblems-problemphyxmini0706-bucketcylindersetup}
  Numerical data and image readouts supplied by the exercise. The source gives the diameter as 4.0 cm, while the image additionally labels the radius as 2.0 cm. The unit-independent statement of release from rest is included for every coherent pair of length and time units
\end{definition}
\begin{proof}
  This is the declared model definition of Matches Problem And Figure Readouts.
\end{proof}
\begin{definition}[Uses Standard Earth Gravity]
  \label{def:physics:phyx-mini-0706:phyxminiproblems-problemphyxmini0706-usesstandardearthgravity}
  \lean{PhyXMiniProblems.ProblemPhyXMini0706.UsesStandardEarthGravity}
  \uses{def:physics:phyx-mini-0706:phyxminiproblems-problemphyxmini0706-accelerationinmeterspersecondsquared, def:physics:phyx-mini-0706:phyxminiproblems-problemphyxmini0706-bucketcylindersetup}
  The standard near-Earth gravitational field used for the numerical multiple-choice evaluation
\end{definition}
\begin{proof}
  This is the declared model definition of Uses Standard Earth Gravity.
\end{proof}
\begin{definition}[Has Physical Parameters]
  \label{def:physics:phyx-mini-0706:phyxminiproblems-problemphyxmini0706-hasphysicalparameters}
  \lean{PhyXMiniProblems.ProblemPhyXMini0706.HasPhysicalParameters}
  \uses{def:physics:phyx-mini-0706:phyxminiproblems-problemphyxmini0706-massreadout, def:physics:phyx-mini-0706:phyxminiproblems-problemphyxmini0706-lengthreadout, def:physics:phyx-mini-0706:phyxminiproblems-problemphyxmini0706-signedlengthreadout, def:physics:phyx-mini-0706:phyxminiproblems-problemphyxmini0706-timereadout, def:physics:phyx-mini-0706:phyxminiproblems-problemphyxmini0706-accelerationreadout, def:physics:phyx-mini-0706:phyxminiproblems-problemphyxmini0706-bucketcylindersetup}
  Positivity and nondegeneracy conditions for the physical parameters
\end{definition}
\begin{proof}
  This is the declared model definition of Has Physical Parameters.
\end{proof}
\begin{definition}[Satisfies Cylinder Geometry]
  \label{def:physics:phyx-mini-0706:phyxminiproblems-problemphyxmini0706-satisfiescylindergeometry}
  \lean{PhyXMiniProblems.ProblemPhyXMini0706.SatisfiesCylinderGeometry}
  \uses{def:physics:phyx-mini-0706:phyxminiproblems-problemphyxmini0706-lengthreadout, def:physics:phyx-mini-0706:phyxminiproblems-problemphyxmini0706-bucketcylindersetup}
  The diameter is twice the radius in every length unit
\end{definition}
\begin{proof}
  This is the declared model definition of Satisfies Cylinder Geometry.
\end{proof}
\begin{definition}[Satisfies Uniform Solid Cylinder Inertia Law]
  \label{def:physics:phyx-mini-0706:phyxminiproblems-problemphyxmini0706-satisfiesuniformsolidcylinderinertialaw}
  \lean{PhyXMiniProblems.ProblemPhyXMini0706.SatisfiesUniformSolidCylinderInertiaLaw}
  \uses{def:physics:phyx-mini-0706:phyxminiproblems-problemphyxmini0706-massreadout, def:physics:phyx-mini-0706:phyxminiproblems-problemphyxmini0706-lengthreadout, def:physics:phyx-mini-0706:phyxminiproblems-problemphyxmini0706-momentofinertiareadout, def:physics:phyx-mini-0706:phyxminiproblems-problemphyxmini0706-bucketcylindersetup}
  Axial moment of inertia of a uniform solid cylinder, I = (1/2) M R²
\end{definition}
\begin{proof}
  This is the declared model definition of Satisfies Uniform Solid Cylinder Inertia Law.
\end{proof}
\begin{definition}[Satisfies Bucket Cylinder Dynamics]
  \label{def:physics:phyx-mini-0706:phyxminiproblems-problemphyxmini0706-satisfiesbucketcylinderdynamics}
  \lean{PhyXMiniProblems.ProblemPhyXMini0706.SatisfiesBucketCylinderDynamics}
  \uses{def:physics:phyx-mini-0706:phyxminiproblems-problemphyxmini0706-massreadout, def:physics:phyx-mini-0706:phyxminiproblems-problemphyxmini0706-lengthreadout, def:physics:phyx-mini-0706:phyxminiproblems-problemphyxmini0706-accelerationreadout, def:physics:phyx-mini-0706:phyxminiproblems-problemphyxmini0706-forcereadout, def:physics:phyx-mini-0706:phyxminiproblems-problemphyxmini0706-momentofinertiareadout, def:physics:phyx-mini-0706:phyxminiproblems-problemphyxmini0706-angularaccelerationreadout, def:physics:phyx-mini-0706:phyxminiproblems-problemphyxmini0706-bucketcylindersetup}
  Governing scalar dynamics, with downward positive for the bucket and the corresponding rotation direction positive for the cylinder: m g - T = m a for the bucket; T R = I α about the central axle; a = R α for a taut, non-slipping string. These laws do not state the requested acceleration or fall time
\end{definition}
\begin{proof}
  This is the declared model definition of Satisfies Bucket Cylinder Dynamics.
\end{proof}
\begin{definition}[Satisfies Constant Acceleration Kinematics]
  \label{def:physics:phyx-mini-0706:phyxminiproblems-problemphyxmini0706-satisfiesconstantaccelerationkinematics}
  \lean{PhyXMiniProblems.ProblemPhyXMini0706.SatisfiesConstantAccelerationKinematics}
  \uses{def:physics:phyx-mini-0706:phyxminiproblems-problemphyxmini0706-signedlengthreadout, def:physics:phyx-mini-0706:phyxminiproblems-problemphyxmini0706-timereadout, def:physics:phyx-mini-0706:phyxminiproblems-problemphyxmini0706-signedspeedreadout, def:physics:phyx-mini-0706:phyxminiproblems-problemphyxmini0706-accelerationreadout, def:physics:phyx-mini-0706:phyxminiproblems-problemphyxmini0706-bucketcylindersetup}
  The vertical coordinate is positive upward. With a constant downward acceleration magnitude a, the endpoint equation is y₁ = y₀ + v₀ t - (1/2) a t². This generic kinematic law constrains the response field fallTime; it does not prescribe its solved value
\end{definition}
\begin{proof}
  This is the declared model definition of Satisfies Constant Acceleration Kinematics.
\end{proof}
\begin{lemma}[bucket Downward Acceleration formula]
  \label{lem:physics:phyx-mini-0706:phyxminiproblems-problemphyxmini0706-bucketdownwardacceleration-formula}
  \lean{PhyXMiniProblems.ProblemPhyXMini0706.bucketDownwardAcceleration_formula}
  \uses{def:physics:phyx-mini-0706:phyxminiproblems-problemphyxmini0706-massreadout, def:physics:phyx-mini-0706:phyxminiproblems-problemphyxmini0706-accelerationreadout, def:physics:phyx-mini-0706:phyxminiproblems-problemphyxmini0706-bucketcylindersetup, def:physics:phyx-mini-0706:phyxminiproblems-problemphyxmini0706-hasphysicalparameters, def:physics:phyx-mini-0706:phyxminiproblems-problemphyxmini0706-satisfiesuniformsolidcylinderinertialaw, def:physics:phyx-mini-0706:phyxminiproblems-problemphyxmini0706-satisfiesbucketcylinderdynamics}
  Eliminating tension and angular acceleration from the three dynamics laws, then using the solid-cylinder inertia law, gives the constant downward acceleration. This is an intermediate consequence, not a premise
\end{lemma}
\begin{proof}
  The conclusion follows from the governing relations and figure readouts named in the statement of bucket Downward Acceleration formula.
\end{proof}
\begin{lemma}[fall Time formula]
  \label{lem:physics:phyx-mini-0706:phyxminiproblems-problemphyxmini0706-falltime-formula}
  \lean{PhyXMiniProblems.ProblemPhyXMini0706.fallTime_formula}
  \uses{def:physics:phyx-mini-0706:phyxminiproblems-problemphyxmini0706-massreadout, def:physics:phyx-mini-0706:phyxminiproblems-problemphyxmini0706-signedlengthreadout, def:physics:phyx-mini-0706:phyxminiproblems-problemphyxmini0706-timereadout, def:physics:phyx-mini-0706:phyxminiproblems-problemphyxmini0706-accelerationreadout, def:physics:phyx-mini-0706:phyxminiproblems-problemphyxmini0706-bucketcylindersetup, def:physics:phyx-mini-0706:phyxminiproblems-problemphyxmini0706-matchesproblemandfigurereadouts, def:physics:phyx-mini-0706:phyxminiproblems-problemphyxmini0706-hasphysicalparameters, def:physics:phyx-mini-0706:phyxminiproblems-problemphyxmini0706-satisfiesuniformsolidcylinderinertialaw, def:physics:phyx-mini-0706:phyxminiproblems-problemphyxmini0706-satisfiesbucketcylinderdynamics, def:physics:phyx-mini-0706:phyxminiproblems-problemphyxmini0706-satisfiesconstantaccelerationkinematics}
  The positive solution of the endpoint kinematics is t = sqrt (2 h (m + M/2) / (m g)), where h = y₀ - y₁. The cylinder radius cancels after the no-slip and solid-cylinder laws are combined
\end{lemma}
\begin{proof}
  The conclusion follows from the governing relations and figure readouts named in the statement of fall Time formula.
\end{proof}
\begin{definition}[Answer Choice]
  \label{def:physics:phyx-mini-0706:phyxminiproblems-problemphyxmini0706-answerchoice}
  \lean{PhyXMiniProblems.ProblemPhyXMini0706.AnswerChoice}
  Answer labels displayed with the exercise
\end{definition}
\begin{proof}
  This is the declared model definition of Answer Choice.
\end{proof}
\begin{definition}[displayed Time In Seconds]
  \label{def:physics:phyx-mini-0706:phyxminiproblems-problemphyxmini0706-displayedtimeinseconds}
  \lean{PhyXMiniProblems.ProblemPhyXMini0706.displayedTimeInSeconds}
  \uses{def:physics:phyx-mini-0706:phyxminiproblems-problemphyxmini0706-answerchoice}
  Times printed beside the answer choices, in seconds
\end{definition}
\begin{proof}
  This is the declared model definition of displayed Time In Seconds.
\end{proof}
\begin{definition}[recorded Dataset Answer]
  \label{def:physics:phyx-mini-0706:phyxminiproblems-problemphyxmini0706-recordeddatasetanswer}
  \lean{PhyXMiniProblems.ProblemPhyXMini0706.recordedDatasetAnswer}
  \uses{def:physics:phyx-mini-0706:phyxminiproblems-problemphyxmini0706-answerchoice}
  Dataset metadata; it is not a premise of the theorem
\end{definition}
\begin{proof}
  This is the declared model definition of recorded Dataset Answer.
\end{proof}
\begin{definition}[Is Unique Closest Displayed Time]
  \label{def:physics:phyx-mini-0706:phyxminiproblems-problemphyxmini0706-isuniqueclosestdisplayedtime}
  \lean{PhyXMiniProblems.ProblemPhyXMini0706.IsUniqueClosestDisplayedTime}
  \uses{def:physics:phyx-mini-0706:phyxminiproblems-problemphyxmini0706-timeinseconds, def:physics:phyx-mini-0706:phyxminiproblems-problemphyxmini0706-bucketcylindersetup, def:physics:phyx-mini-0706:phyxminiproblems-problemphyxmini0706-answerchoice, def:physics:phyx-mini-0706:phyxminiproblems-problemphyxmini0706-displayedtimeinseconds}
  The physical fall time is strictly closer to one displayed value than to every other displayed value
\end{definition}
\begin{proof}
  This is the declared model definition of Is Unique Closest Displayed Time.
\end{proof}
% --- Archon declaration topology end ---
% --- Archon physics formalization source end ---
